\section{Dataset}
\label{sec:dataset}

% Cross-reference: full Datasheets-for-Datasets card lives at DATA_CARD.md
% in the repository root. This section summarises only what is necessary
% to evaluate the benchmark.

\subsection{Substrate and provenance}

All tasks are grounded in the publicly redistributable CDISC SDTM/ADaM Pilot
Project (CC-BY-4.0), pinned via Git submodule at commit
\texttt{658fcc05506b169a27dee6e2c3a1ccdaaf64a716}. Each task's
\texttt{provenance} block records its derivation script and the substrate
commit, making generation end-to-end auditable. A SHA-256 manifest of 272
fixture files is committed alongside the tasks.

\paragraph{Vendored fixture subset.} For grounding-conditioned baselines
(\S\ref{sec:baselines}), a 26-file, 682\,KB subset of the substrate is
vendored verbatim under \texttt{fixtures/}. The subset is the
intersection of the 67 fixture base paths referenced by
\texttt{inputs.fixtures} across the 250 task specifications and the
files that exist in the v1.0.0 substrate; the remaining 42 references
(15 SDTM domain-creator R programs and 27 TLF mock templates with
hyphen/underscore naming drift) are unauthored upstream and are tracked
as v0.1.1 work (\S\ref{sec:limitations}). Vendoring is verbatim with no
modification; substrate \texttt{LICENSE} and \texttt{NOTICE.md} accompany
the files.

\subsection{Composition}

ClinProg-Bench contains \textbf{250} tasks distributed across five categories
(Table~\ref{tab:taxonomy}). Each task is a validated \texttt{TaskSpec} JSON
file bundled with fixture files, a markdown specification, and a gold output
under \texttt{gold/<task\_id>/}.

\begin{table}[h]
\centering
\small
\caption{Task taxonomy. Output kinds and scorers are fixed per
sub-category to support deterministic evaluation.}
\label{tab:taxonomy}
\begin{tabular}{llrll}
\toprule
Category & Sub-categories & N & Output kind & Scorer \\
\midrule
T1 Code Generation & SDTM (22), ADaM (18), TLF (20)         & 60 & program & \texttt{codegen} \\
T2 Code Review     & program (25), validation (25)          & 50 & JSON    & \texttt{log\_review} \\
T3 Spec Extraction & SDTM (21), ADaM (18), TLF (11)         & 50 & JSON    & \texttt{spec\_extract} \\
T4 Documentation   & TLF (16), dataset card (13), synopsis (11) & 40 & text & \texttt{doc} \\
T5 Debugging       & SDTM (17), ADaM (17), TLF (16)         & 50 & patch   & \texttt{debug} \\
\midrule
Total              &                                        & 250 & & \\
\bottomrule
\end{tabular}
\end{table}

\subsection{Quality assurance}

\paragraph{Schema validation.} Every task is parsed by the strict pydantic
schema in \texttt{src/clinprog\_bench/schema.py}; \texttt{task\_id} is fixed
to the regex \texttt{T[1-5]\textbackslash.[1-9][0-9]?\textbackslash.[a-z0-9\_]+\textbackslash.[0-9]\{3\}}
and \texttt{seed\_commit} to a 40-char SHA-1. Validation runs on every CI
push (\texttt{.github/workflows/eval.yml}, job \texttt{validate}).

\paragraph{Leakage audit.} A two-channel audit (\texttt{scripts/leakage\_audit.py})
checks (a) SHA-256 byte overlap between gold outputs and the substrate, and
(b) 13-gram Jaccard similarity between task prompts and substrate text. Both
channels report \textbf{0 hits} on the v0.1.0 release; the report is
committed at \texttt{docs/leakage-audit.md}.

\paragraph{Dual-reviewer audit.} A stratified random subset (seed
\texttt{20260426}, 4 tasks per category, 20 total) is independently rated by
two domain experts using a 5-point rubric. The acceptance threshold is Cohen's
$\kappa \geq 0.8$. Materials are shipped under \texttt{review\_package/} and
the audit log template is at \texttt{docs/audit-log.md}.

\subsection{Licensing and release}

The benchmark toolkit (Python package, scorers, runner contract) is released
under \textbf{Apache-2.0}. The task corpus (JSON, gold outputs, fixtures) is
released under \textbf{CC-BY-4.0}, inheriting the substrate license. A
versioned snapshot is archived on Zenodo
(DOI withheld for double-blind review).
A full Gebru-style data card lives at \texttt{DATA\_CARD.md} in the repository
root.
